% Volume V: Admissibility and Continuation
% Part IX. A General Theory of Institutional State Change
% Chapter 48: Pop, Refuse, Bind, Collapse
% Spine: proof-spines/ch048-spine.md

\chapter{Pop, Refuse, Bind, Collapse}
\label{ch:ch048}

The Spherepop primitives are translated here into legal and epistemic operations. Let \(\mathsf P,\mathsf R,\mathsf B,\mathsf C\) denote \textbf{Pop}, \textbf{Refuse}, \textbf{Bind}, and \textbf{Collapse}. Pop introduces a distinction into active consideration. Refuse blocks a proposed transition. Bind joins evidence, claim, person, rule, or responsibility under a relation. Collapse converts an unresolved field of alternatives into an operative determination.

The chapter then develops an operational algebra and asks which compositions commute:
\[
\mathsf C\circ\mathsf B
\stackrel{?}{=}
\mathsf B\circ\mathsf C.
\]
\ProposedTheorem\ the chapter's load-bearing claim is that these compositions generally fail to commute in legal settings.
In most legal settings they do not. Binding evidence to a person before interpretive collapse differs from collapsing a narrative and only then assigning the person to it. The noncommutation matters because institutional consequence is order-sensitive: the same primitives can yield different public states depending on when refusal, binding, or collapse occurs.
